\section{Extensibility, Inspectability and Expressiveness}
\label{sec:features}

The features seen in Section~\ref{sec:motivation} form the basis of \dsl{} as a library. However, many animation libraries such as GreenSock Animation Platform provide a much more extensive list of features.

To add support for similar features, we design \dsl{} with extensibility and inspectability in mind. This means that \dsl{} can be extended with custom operations and information can be derived from inspecting specified animations. To be able to support arbitrary expressiveness in combination with those features, we also need to stress the possibility to extend \dsl{} with custom combinators.

\subsection{Extensibility}
\label{sec:customop}

The \hs{linearTo} operation and the \hs{sequential} and \hs{parallel} combinators form the basis for expressing a variety of animations. However, there are different situations which require different primitives to express our desired animation. For example, GSAP provides a primitive to morph one shape into another.

An example in our use case is the \hs{checkIcon} animation, which is part of the \hs{markAsDone} animation, where it is required to set the color of the checkmark to a specific value. We can define an additional \hs{set} operation and embed it inside a \dsl{} animation. In this animation we use Haskell's \hs{do}-notation to specify sequential animations.

\begin{spec}
checkIcon = do
  linearTo (checkmark . scale) (For 0.05) (To 0)
  set (checkmark . color) completeGreen
  (linearTo (circle . extra) (For 0.2) (To 360))
    `parallel`
    (linearTo (checkmark . scale) (For 0.2) (To 1.2))
  linearTo (checkmark . scale) (For 0.05) (To 1)
\end{spec}

%\begin{figure}[h]
%\centering
%\includegraphics[width=\figscale\textwidth]{pictures/todo}
%\caption{The \hs{checkIcon} animation.}
%\label{fig:}
%\end{figure}

\subsection{Inspectability}

\dsl{} is inspectable, meaning that we can derive properties of expressed computations by \emph{inspecting them rather than running them}. For example, if we want to know the duration of \hs{menuIntro} without actually running it and keeping track of the time. We can do this by using a predefined \hs{duration} function, which calculates the duration by inspecting the animation. This gives a duration of \hs{0.5} seconds, which is indeed the duration of two \hs{0.5} second animations in parallel.

\begin{spec}
menuIntroDuration = duration menuIntro
-- menuIntroDuration = 0.5 
\end{spec}

Of course, it is not possible to inspect every animation. In the following situation we have a custom operation \hs{get}, the dual of \hs{set} in the previous section, returning a \hs{Float}. If the result of this value is used as the duration parameter of an animation, then we cannot know upfront how long this animation will last. Requesting to calculate the duration of such an animation results in a type error.

\begin{spec}
complicatedAnimation = do
  v <- get
  linearTo lens (For v) (To 10)

complicatedAnimationDuration = duration complicatedAnimation
-- type error
\end{spec}

Calculating a duration is mostly a stepping stone towards supporting other interesting features. One such example is sequentially composing animations with a relative offset. For example, to compose a first animation \hs{anim1} with a second animation \hs{anim2} which starts 0.5 seconds before the end of \hs{anim1}.

\begin{spec}
relSeqAnim = relSequential anim1 anim2 (-0.5)
\end{spec}

\subsection{Expressiveness}
\label{sec:customcomb}

Similarly to providing custom operations, \dsl{} also supports custom combinators. For typical DSLs this is not a requirement since working with monadic computations provides the combinators \hs{>>=} and \hs{return} which are suitable for the required use cases. However, since \dsl{} has the additional requirement of being inspectable, the \hs{>>=} combinator ends up being a liability because it only provides a very limited amount of inspectability.

In the \hs{onlyDone} animation, we show all \emph{done} to-do items, while hiding all \emph{not done} to-do items. This could be implemented by first showing all items with the \hs{showAll} animation, since an item might have been hidden by a previous action, and then hiding all not done items with the \hs{hideNotDone} animation. The definition for this is given below, while the implementation of \hs{showAll} and \hs{hideNotDone} is omitted.

\begin{spec}
onlyDoneNaive = do showAll ; hideNotDone
\end{spec}

However, this animation is a bit naive since it executes the \hs{showAll} animation regardless of whether there are any hidden done items that actually need to appear. Instead, the intended animation is to only make done items appear when some were hidden. This is done by first checking whether there are any done items and based on that we play the naive version of \hs{onlyDone}, otherwise we just hide the not done items.

\begin{spec}
onlyDone = do
  cond <- doneItemsGt0    -- check if more than 0 'done' items
  if cond then onlyDoneNaive else hideNotDone
\end{spec}

However, if we also want to inspect this computation the previous formulation is problematic since it is specified using a construction which is too general, we explain this in more detail in Section~\ref{sec:interaction}. Instead, we can define a custom combinator \hs{ifThenElse} which captures the form of animation we want to express. With this combinator we can express \hs{onlyDone} in an inspectable manner.

\begin{spec}
onlyDone = ifThenElse doneItemsGt0 onlyDoneNaive hideNotDone
\end{spec}

For this new combinator, we can define custom ways to inspect it. Since each branch might have a different duration, it does not make sense to extract a duration from it. Instead we might want to extract the \emph{maximum} duration of the animation.

\begin{spec}
onlyDoneMaxDuration = maxDuration onlyDone
-- onlyDoneMaxDuration = 1
\end{spec}

Section~\ref{sec:motivation} and \ref{sec:features} gave a look and feel of the features of \dsl{}. In the following sections, we delve deeper into the internals of the implementation.
